\subsection{Drift by family and control up to the first exit}\label{macro:control}

\begin{lema}[Conditional loss of a pattern]\label{macro:deriva}
Suppose the parent satisfies the thresholds of Proposition~\ref{macro:maximo} and $k\ge\delta n$.
For a set $\mathcal L$ of boundedly many distinct available lines, let $H^*(\mathcal L)$ be the sum
of one for each row or column and $\ell_f^*$ for each diagonal or antidiagonal $f$. In one step,
\[
 P^{\rm B}(\text{some line in }\mathcal L\text{ becomes occupied}\mid\mathcal F_s)
 =\frac{H^*(\mathcal L)}{k}
       +O_\delta\!\left(\frac{\zeta+k^{-1}}{k}\right).
\]
The fiber rates are those of Definition~\ref{macro:fibras}, and the rate for the rectangle term $T$
is
\[
 H_T^*=4+\sum_{m\in T}\bigl(\ell_{D(e_m)}^*+\ell_{A(e_m)}^*\bigr).
\]
\end{lema}
\begin{proof}
Lemma~\ref{macro:casilla} gives the sum of the individual probabilities. Replacing diagonal masses by
their targets incurs the error in \eqref{macro:predictor} divided by $k$ and multiplied by the number
of lines. The intersection of any two distinct lines has probability at most
$\max_e p_e/k=O_\delta(k^{-2})$. The first two inclusion--exclusion inequalities bound the error for
the union by the sum of these intersection probabilities.

For an available tracked line $\lambda$, write $Y_{e|\lambda}=\prod_{f\ni e,\,f\ne\lambda}X_f$.
If $Y_{e|\lambda}=1$, its contribution decreases when any of the other three lines is hit while
$\lambda$ remains unoccupied. Excluding simultaneous occupation of $\lambda$ and one of the three
lines changes the probability by $O_\delta(k^{-2})$. If $\lambda$ becomes occupied, the freezing
rule preserves the observable. Thus, writing $\nabla_s$ for the one-step increment, we have
\[
 \E_{\rm B}[\nabla_s F_{\lambda,I}\mid\mathcal F_s]
 =-\frac1{kn}\sum_{e\in\lambda_I}Y_{e|\lambda}H_\lambda^*(e)
   +O_\delta\!\left(\frac{\zeta+k^{-1}}{k}F_{\lambda,I}\right).
\]
By \eqref{macro:derivadas}, the time derivative of each fiber density in the table is $-H^*_\lambda/x$
times the density itself. Without freezing, the increment of its target is the integral of this
derivative divided by $n$, with Taylor error $O_\delta(n^{-2}|\lambda_I|)$, where $|\lambda_I|$ is
the length of its parameter interval. The probability of freezing is $O(1/k)$, and the ordinary
increment of the target is $O_\delta(|\lambda_I|/n)$. Freezing the target at the step when the line
becomes occupied therefore adds $O_\delta(|\lambda_I|/(kn))$ to its conditional drift. For a fiber
that is already frozen, both drifts are zero.

For row or column masses we have exactly $\E_{\rm B}[\nabla_s N_I\mid\mathcal F_s]=-N_I/k$, and the
target decreases by $|I|/n$. For the rectangle, we apply the union formula to the four row and column
lines and the $2|T|$ diagonal lines. These are distinct for ordinary quadruples. If $Y_{T}$ denotes
the product of their indicators, then
\[
 \E_{\rm B}[\nabla_s P_{\Theta,T}\mid\mathcal F_s]
 =-\frac1{kn^4}\sum_{(i,i',j,j')\in\Theta}^{\rm ord}Y_T H_T^*
   +O_\delta\!\left(\frac{\zeta+k^{-1}}{k}P_{\Theta,T}\right).
\]
When $Y_T=0$, its contribution remains zero. The derivative of the target density is exactly
$-H_T^*/x$ times the density itself. The second time derivative is uniformly bounded for $x\ge\delta$,
so the integrated Taylor error is $O_\delta(|\Theta|/n^2)$.
\end{proof}

\begin{proposicion}[Stopping bound under the auxiliary law]\label{macro:parada}
Let $M=\lceil\delta n\rceil$. For every sufficiently small radius $\varrho_{\rm mac}>0$, we can fix a
partition $\mathcal I$ of mesh $a_{\rm mac}$ and constants
$c_{\delta,\varrho_{\rm mac}},C_{\delta,\varrho_{\rm mac}}>0$ with the following property.
Let $\mathcal H_M$ be the event that the discrepancy reaches or exceeds $\varrho_{\rm mac}$ at some
board of size ${k\ge M}$ in the history, including the board of size $M$. For sufficiently large $n$,
\[
 P^{\rm B}(\mathcal H_M)\le C_{\delta,\varrho_{\rm mac}}\exp(-c_{\delta,\varrho_{\rm mac}}n).
\]
A fixed observation partition may be monitored on the same histories by taking $\mathcal I$ to
refine it. Outside $\mathcal H_M$, its discrepancy stays below the same radius through size $M$.
\end{proposicion}
\begin{proof}
Let $\tau$ be the index of the first board of size at least $M$ in the history whose discrepancy
reaches or exceeds $\varrho_{\rm mac}$, and set $\tau=n-M$ if there is no such board. This is a
stopping time with $\tau\le n-M$. Every board of index less than $\tau$ has discrepancy less than
$\varrho_{\rm mac}$ and, by Proposition~\ref{macro:maximo}, has a Bethe point. Thus the history
stopped at $\tau$ is defined under $P^{\rm B}$, and $\mathcal H_M=\{\Delta(\tau)\ge\varrho_{\rm mac}\}$
depends only on that history. We stop the discrepancies and the centered increments of each
observable minus its target at $\tau$. We include the increment of the step reaching $\tau$, since
the preceding board is still within the radius. Let $\mathcal M^*$ be the sum of the maxima over
time of the absolute values of these martingales for row and column intervals and for boxes and
subsets, plus the maximum over tracked lines of the sum of these absolute maxima for their intervals.

In a box or interval, we subtract from the empirical and target drifts the contributions of the same
rate $H^*$, evaluated at a point in that region. The constant-rate term is at most $C_\delta/n$ times
the absolute mass discrepancy in the region. The oscillation of $H^*$ is $O_\delta(a_{\rm mac})$ by
spatial regularity and multiplies the sum of the empirical and target masses. Over all subfibers of a
line these masses sum to at most two. For each of the sixteen types they likewise sum to at most two
over all boxes. In the Taylor errors, the lengths or volumes sum to at most one per family or type.
These facts yield, with a constant $C_\delta$ independent of the number of intervals,
\begin{equation}\label{macro:gronwall}
 \Delta(s\wedge\tau)\le \Delta(0)+\mathcal M^*
 +\frac{C_\delta}{n}\sum_{j<s\wedge\tau}\Delta(j)
 +C_\delta(a_{\rm mac}+C_{\mathcal I}/n),\qquad 0\le s\le n-M .
\end{equation}
For the fibers, we first sum the absolute errors over their intervals and then take the maximum over
lines. Once a fiber freezes, its increments vanish and its discrepancy at freezing remains in
$\Delta_F$. In estimating the drift,
Lemma~\ref{macro:deriva} uses only $\zeta$, which is linear in $\Delta_0+\Delta_F$. The exponent
$1/4$ in Lemma~\ref{macro:brecha} does not affect the coefficient in \eqref{macro:gronwall}.

There are $O_{\mathcal I}(n)$ scalar martingales. For a fiber that remains available, each newly
occupied line affects at most one square of the fiber, since distinct lines intersect at most once.
The step that occupies the tracked line freezes the entire observation for that line. For a
rectangle, requiring one of the occupied lines to pass through a given position in the rectangle
imposes a nontrivial linear relation and allows $O(n^3)$ quadruples. Each increment of a normalized
mass, including the target increment, therefore has absolute value at most $C_\delta/n$. The same
bound, with a different constant, holds after centering and stopping the increments.
